\documentclass[11pt]{article}
\usepackage{icmlko}
\begin{document}

\title{해는 같고 값만 다르다 --- 큰 도메인은 누구든 남에게 $+0.003$ 을 물린다}
\author{월드모델 실험실}
\date{노트 364 · 2026년 8월 1일}
\maketitle

\begin{abstract}
노트 363이 복도 위험 지도를 그리고 예측을 미리 적었다 --- \emph{복도 위험이
해를 예측하면, 위험 0.1922 인 세계애니의 학습을 빼는 쪽이 위험 0.0539 인
만화를 빼는 쪽보다 남에게 더 큰 이득을 줘야 한다.} 대상과 가중이 같도록
\textbf{만화와 세계애니를 뺀 아홉 도메인}으로만 채점했다.
\textbf{예측이 깨졌다} --- 만화 $+0.0035$(SD 0.0084, 8/12) · 세계애니
$+0.0030$(SD 0.0076, 9/12)로 \textbf{같다}. 위험은 3.6배, 행 수는 1.5배
차이인데 남에게 준 해는 같다. \textbf{그러니 해는 복도 위험도 행 수도 아니고,
큰 도메인이면 대체로 같은 값($\approx +0.003$)을 남에게 물린다.}
둘을 가르는 것은 \emph{값을 치르느냐}뿐이다 --- 세계애니는 제 점수를
$-0.1474$(0/12) 잃고 판의 11\%를 벌어 오지만 만화는 아무것도 안 벌어 온다.
\textbf{만화 문제는 처음부터 회계였다.} 그리고 셋째 사실이 하나 더 나왔다:
\textbf{팝업은 둘 다에서 잃는다}($-0.0531$ · $\mathbf{-0.0993}$, 0/12).
학습 24행짜리 도메인이 남의 학습에 제일 크게 기대고 있다 --- 이 실험실의
결정 넷을 연속으로 막은 것이 왜 팝업이었는지가 여기서 설명된다.
\end{abstract}

\section{공정한 검사}

노트 363이 지도를 그리고 만화가 11위인 것을 보고 물었다 --- 그러면 왜
만화를 빼야 판이 오르나. 답 후보가 둘이었다: 위험이 실제로 해를 만드는데
지도가 잘못 됐거나, 위험이 아니라 \emph{회계}이거나.

\textbf{대상 집합을 고정해 가른다.} 세 팔 모두 \emph{만화와 세계애니를 뺀
아홉 도메인}으로만 채점한다(분모 2,423). 그러면 각 팔에서 빠지는 도메인
자신의 점수가 판에 안 들어가므로 \textbf{남에게 준 해}만 남는다.

\begin{center}\includegraphics[width=\linewidth]{figs/same.pdf}\end{center}

\begin{center}
\begin{tabular}{lrrrrr}
\hline
학습을 뺀 도메인 & 행 & 복도 위험 & 아홉 판 차 & SD & 부호 \\
\hline
만화     & 1,783 & 0.0539 & $+0.0035$ & 0.0084 & 8/12 \\
세계애니 & 2,648 & \textbf{0.1922} & $+0.0030$ & 0.0076 & 9/12 \\
\hline
\end{tabular}
\end{center}

\textbf{같다.} 미리 적은 예측은 세계애니 쪽이 3.6배 커야 한다는 것이었다.
행 수로도 1.5배 차이인데 그것도 안 나타난다.

\textbf{그러니 남에게 주는 해는 복도 위험의 함수가 아니다.} 이 판에서
학습의 17\textasciitilde{}25\% 를 차지하는 도메인은, 어느 쪽이든 나머지에게
대략 $+0.003$ 어치의 자리를 빼앗고 있다. 노트 362의 문턱(1.7\%\textasciitilde
{}8.5\%)과 합치면 그림이 단순해진다 --- \textbf{몫이 문턱을 넘으면 값이
붙고, 넘고 나면 크기를 별로 안 탄다.}

\section{그러면 왜 만화만 빼는 게 이득인가 --- 회계}

\begin{center}
\begin{tabular}{lrrr}
\hline
 & 남에게 준 해 & 제 점수 변화 & 판 가중 \\
\hline
만화     & $+0.0035$ & --- (채점 안 됨) & \textbf{0\%} \\
세계애니 & $+0.0030$ & $\mathbf{-0.1474}$ (0/12) & 11.0\% \\
\hline
\end{tabular}
\end{center}

세계애니 학습을 빼면 남이 $+0.0030$ 얻는 대신 세계애니 자신이 $0.1474$ 를
잃는다. 가중 11\% 를 곱하면 $-0.0162$ --- 전체 판으로는 크게 손해다.
만화는 그 두 번째 칸이 \textbf{비어 있다.} 유보 6행이라 \texttt{evaluate}
문턱 20을 못 넘고, 그래서 잃을 점수가 없다.

\textbf{노트 356부터 362까지 넉 편이 쫓던 것이 이것이다.} 만화가 특별히
나쁜 자료여서가 아니다 --- \emph{똑같이 나쁜} 자료인데 혼자만 값을 안 낸다.
노트 363이 ``해가 아니라 회계''라고 적은 것이 그대로 확인됐다.

\section{덤 --- 팝업은 왜 늘 반대편인가}

\begin{center}
\begin{tabular}{lrrrr}
\hline
도메인 & 학습 & 만화 뺌 & 세계애니 뺌 & \\
\hline
\textbf{팝업} & \textbf{24} & $-0.0531$ (2/12) & $\mathbf{-0.0993}$ & \textbf{(0/12)} \\
아이돌 & 122 & $+0.0552$ (10/12) & $+0.0632$ & (10/12) \\
게임   & 259 & $-0.0009$ & $+0.0118$ & (10/12) \\
펀딩   & 320 & $+0.0020$ & $+0.0263$ & (10/12) \\
웹툰   & 2,106 & $+0.0097$ & $+0.0017$ & \\
\hline
\end{tabular}
\end{center}

\textbf{팝업은 큰 도메인의 학습이 빠질 때마다 잃는다} --- 세계애니를 빼면
$-0.0993$ 에 0/12 다. 학습 24행짜리 도메인은 제 힘으로 배우는 게 거의
없고 공유 적합이 남의 행에서 배운 것을 빌려 쓴다. \textbf{노트 252의
``풀링은 타협이 아니라 이득''이 제일 센 형태로 나타난 자리다.}

그리고 이것이 \textbf{이 실험실의 결정 넷을 연속으로 막은 이유}다 ---
351(팝업 올라서 채택) · 356(팝업 내려서 안 뺌) · 362(팝업 내려서 안 넣음) ·
그리고 여기. 판을 올리는 처치는 대개 \emph{학습을 줄이는} 처치이고,
학습을 줄이면 팝업이 제일 먼저 다친다. \textbf{대표 수치와 판이 구조적으로
반대 방향을 본다.}

\section{규약에 더하는 줄}

\begin{itemize}
\item \textbf{``남에게 준 해''는 대상 집합을 고정해서 잰다.} 도메인을
빼면서 그 도메인 점수를 판에 남겨 두면 자기 손해와 남의 이득이 섞인다.
\item \textbf{큰 도메인의 해는 크기를 별로 안 탄다.} 몫이 문턱(1.7\%
\textasciitilde{}8.5\%)을 넘으면 $\approx +0.003$ 이고, 그 위로는 위험도
행 수도 예측력이 없다.
\item \textbf{학습을 줄이는 처치는 제일 얇은 도메인부터 다친다.} 팝업
학습 24행이 이 판의 구조적 제약이다.
\end{itemize}

\section{남은 것}

\begin{center}
\begin{tabular}{ll}
\hline
갈래 & 지금 아는 것 \\
\hline
\textbf{만화 2026 스무 건} & 회계를 닫는다. 남에게 주는 $+0.0035$ 를 \\
 & 제 점수로 회수한다. \textbf{다음 순번} \\
팝업 실측 41행 & 노트 360. 결정 넷을 다 팝업이 내렸다 \\
문턱이 정확히 어디인가 & 1.7\%\textasciitilde{}8.5\%(노트 362). 4\% 팔 하나면 좁아진다 \\
아이돌과 팝업이 왜 반대인가 & 학습 122 대 24. 풀링이 돕다가 해치는 \\
 & 자리가 그 사이에 있나 \\
\hline
\end{tabular}
\end{center}

\end{document}
